\section{Diagênese e litificação}

\begin{definicao}
\textbf{Diagênese} é o conjunto de processos físicos, químicos e biológicos que modifica o sedimento após a deposição e durante o soterramento, sob condições inferiores às do metamorfismo. A \textbf{litificação} corresponde à transformação do sedimento inconsolidado em rocha sedimentar.
\end{definicao}

\subsection{Compactação}

A compactação reduz o volume de poros por rearranjo, deformação ou fraturamento dos grãos. Em sedimentos lamosos, a expulsão de água intersticial pode produzir grande redução de volume. A compactação química pode ocorrer nos contatos submetidos a maior tensão, favorecendo dissolução por pressão.

\subsection{Cimentação}

A cimentação resulta da precipitação de minerais nos poros. Cimentos comuns incluem sílica, calcita, dolomita e óxidos ou hidróxidos de ferro. O cimento une os grãos e reduz a porosidade, embora dissoluções posteriores possam gerar porosidade secundária.

\subsection{Dissolução, recristalização e substituição}

Durante a diagênese, grãos ou cimentos instáveis podem ser dissolvidos. A \textbf{recristalização} altera o tamanho ou a forma dos cristais sem necessariamente mudar a composição mineral. Na \textbf{substituição}, um mineral é dissolvido enquanto outro precipita em seu lugar, podendo preservar parcialmente a forma original.

\subsection{Porosidade e permeabilidade}

A \textbf{porosidade} é a proporção de espaços vazios em relação ao volume total da rocha. A \textbf{permeabilidade} expressa a facilidade com que um fluido atravessa uma rede de poros conectados. Uma rocha pode apresentar porosidade elevada e baixa permeabilidade quando os poros são muito pequenos ou pouco conectados.

\subsection{Limite com o metamorfismo}

Não existe uma fronteira única e universal entre diagênese e metamorfismo. A distinção é feita a partir da associação mineral, da temperatura, da pressão e do contexto geológico. Por isso, a diagênese deve ser entendida como uma evolução progressiva, e não como uma etapa instantânea \parencite{nichols2009}.
